\begin{recipe}
[
    preparationtime = {\unit[20]{min}},
    portion = {\portion{2}},
    calory = {260 kcal | 5g Protein | 32g KH | 11g Fett},
    rating = {4},
    source = {Nico Skerra},
    price = {0,74€}
]
{Zwiebelsuppe}

\introduction{%
    Dieses Rezepte stammt nicht von mir direkt, sondern habe ich von einem Freund erhalten.
}

\ingredients{%
    2 & Große Zwiebeln\index{Gemüse \& Pilze!Zwiebeln}\\
    1 EL & Mehl\index{Getreide, Pasta \& Brot!Weizenmehl}\\
    50 ml & Weißwein\index{Getränke!Weißwein}\\
    1 EL & Tomatenmark\index{Vorrat \& Konserven!Tomatenmark}\\
    400 ml & Gemüsebrühe\index{Vorrat \& Konserven!Gemüsebrühe}\\
    10 g & Butter oder Olivenöl\index{Milchprodukte \& Alternativen!Butter}\index{Öle, Saucen \& Würzmittel!Olivenöl}\\
    & Salz\index{Kräuter \& Gewürze!Salz}\\
    & Pfeffer\index{Kräuter \& Gewürze!Pfeffer}\\
    (optional) & Petersilie\index{Kräuter \& Gewürze!Petersilie}\\
    (optional) & 1 EL Zucker\index{Backzutaten \& Süßes!Zucker}
}

\preparation{%
    \step%
    Butter oder Öl erhitzen.\\

    \step%
    Zwiebeln in Halbringe schneiden und langsam anschwitzen, bis sie leicht bräunen und Bratensatz entsteht, hier eventuell den Zucker zufügen.

    \step%
    Tomatenmark unterrühren, dann Mehl einarbeiten.\\

    \step%
    Mit Weißwein ablöschen und etwas einkochen lassen.\\

    \step%
    Mit Gemüsebrühe aufgießen und einige Minuten köcheln.
    
    \step%
    Mit Salz, Pfeffer und optional Petersilie abschmecken.
}

\suggestion[Upgrade]{%
    Mit geröstetem Brot und etwas Käse wird daraus eine schnelle Gratinsuppe.
}

\hint{%
    Zwiebeln nicht zu schnell bräunen, mittlere Hitze gibt mehr Süße und weniger Bitterkeit.
}

\end{recipe}